\section{Security Behavior}
\label{sec:security}

This section states what \texttt{fsb} does to protect the user, as behavior that can
be observed.  The reasoning, and the ways it was tested, are in the design
specification and the security policy.

\subsection{Reach}

\begin{req}{SEC-1}
\texttt{fsb} listens only on the loopback address \code{127.0.0.1}.  No flag,
setting or environment variable changes this, and it is not reachable from another
machine.
\end{req}

\begin{req}{SEC-2}
A request is refused (status 403) unless its \code{Host} header is exactly
\code{127.0.0.1:PORT} or \code{localhost:PORT}, so that a web page on another name
that resolves to the loopback address (DNS rebinding) cannot use it.
\end{req}

\begin{req}{SEC-3}
A request is refused (status 403) if it carries an \code{Origin} or \code{Referer}
from any other origin, if its \code{Sec-Fetch-Site} says it was started by another
origin or another site (including a page on another port of the same host), or if
its \code{Origin} is the opaque value \code{null}.  No cross-origin access is
permitted.
\end{req}

\begin{req}{SEC-4}
Only the methods \code{GET} and \code{HEAD} are answered; anything else gets status
405.  \texttt{fsb} never modifies, creates, renames, moves or deletes anything that
it serves.  The only files it writes are its own, in \code{\~{}/.config/fsb/}: the
rule files of \code{--init} (\rid{CLI-7}), which never overwrites one, the remembered
port (\rid{CFG-5}), and the process-number and log files of a background instance
(\rid{CFG-6}).  Before writing one of them \texttt{fsb} checks that it is an
ordinary file with a single name that belongs to the user, and opens it without
blocking and without emptying it first, so that a link, a pipe or another user's file
placed there cannot make \texttt{fsb} overwrite something else.
\end{req}

\subsection{Credential}

\begin{req}{SEC-5}
Access requires the session, which is established as follows.  At each start
\texttt{fsb} generates a random launch token (at least 128 bits), a random session
value, and a random path prefix (at least 128 bits).  The launch URL carries the
token.  The first request that presents the correct token is answered with a
redirect that removes the token from the address bar, and sets the session cookie.
The token then no longer works, so a copy of the URL left in shell or browser
history is useless.
\end{req}

\begin{req}{SEC-6}
The session cookie is \code{HttpOnly} (scripts cannot read it) and
\code{SameSite=Strict}.  Its path is the random prefix.
\end{req}

\begin{req}{SEC-7}
Browsers scope cookies by host and never by port, so a plain cookie would be sent to
every other web server on \code{127.0.0.1} that the same browser visits.  For that
reason every URL of \texttt{fsb} lies under the random prefix (\code{/PREFIX/}), and
a request whose path is not under it is refused (status 403) like an unauthenticated
request, whatever cookie it carries.  The interface forms every request relative to
its own location, so the cookie accompanies exactly the requests that need it.
\end{req}

\begin{req}{SEC-8}
A request without a valid session is refused with status 403 and the same body,
whatever it asked for.  The launch token is never written to the log.
\end{req}

\subsection{Responses}

\begin{req}{SEC-9}
Every response carries \code{X-Content-Type-Options: nosniff},
\code{Cache-Control: no-store}, \code{Referrer-Policy: no-referrer},
\code{Cross-Origin-Resource-Policy: same-origin}, and, for the interface itself,
\code{X-Frame-Options: DENY} and a content security policy that permits only
resources of its own origin, forbids framing, and forbids form submission and base
URLs.  The interface cannot be framed by another page.
\end{req}

\begin{req}{SEC-10}
A file is never displayed from the interface's own origin except as one of the
previews of Section~\ref{sec:preview}.  Downloading (\rid{API-4}) always answers
with \code{Content-Type: application/octet-stream} and
\code{Content-Disposition: attachment}, so a browser saves the file and never
renders it.
\end{req}

\begin{req}{SEC-11}
Inline responses other than the interface itself are limited to: raster images
whose type is confirmed by their bytes (PNG, JPEG, GIF, WebP), which are sent with a
policy that sandboxes them and allows nothing else; PDF documents whose bytes begin
with \code{\%PDF-}, which may be framed only by the interface's own pages; PNG
pictures that macOS Quick Look draws of iWork and Office documents (\rid{PRV-12}),
sent like the raster images; and the
fixed page that renders Markdown, which runs in a sandboxed frame that has no
access to the session or the interface, may run only its own fixed script, and may
load nothing but images given to it as data.  File names are never trusted.
\end{req}

\begin{req}{SEC-12}
\texttt{fsb} makes no outbound network connection, sends no telemetry, and does not
update itself.  Previewing a file never loads anything remote (\rid{PRV-7}).
\end{req}

\begin{req}{SEC-15}
\texttt{fsb} runs no other program on the user's files except Apple's
\code{qlmanage}, to draw a Quick Look picture (\rid{PRV-12}), and only on macOS.  It is never given
the user's own file: it is given a private copy of a document that every other
endpoint would serve, which \texttt{fsb} read after the same checks, as a single
argument and never through a shell.  The copy carries no extended attributes or
resource fork.  It is stopped after 5 seconds, at most two run at once, and one
whose request is abandoned is stopped.  Copying the document for it is limited to
20 seconds, cooperatively (\rid{PRV-12}), and also stops when the request is
abandoned.
\end{req}

\begin{req}{SEC-13}
\texttt{fsb} never causes a cloud-only file or folder to be downloaded.  The details, the
hover peek, every preview, and the archive listing check whether a \emph{file} is stored
only in the cloud, and listing and search check whether a \emph{folder} is (a folder can be
a cloud-only placeholder in its own right, not only the files inside it); each reports
``not downloaded'' (status 409) instead of reading it.  The one exception is an explicit
download by the user of a cloud-only file (\rid{API-4}), which is the user's request for
it; there is no equivalent for a folder, since fsb never downloads or extracts one.
The same holds for the hard-link index walk (\rid{RUL-6}) and the copy of a package
made for Quick Look (\rid{PRV-12}): neither enters a cloud-only folder, and a package
with a cloud-only part is refused.
\end{req}

\begin{req}{SEC-14}
A path that is refused for any reason other than a permission error on a
non-denied path is answered as not found (status 404) with one and the same body.  In
particular a denied path, a path outside the roots, a path that does not exist, and a
path that is not a regular file or folder cannot be told apart.  Timing is not
equalized: a denied path can be answered slightly faster than a missing one.
\end{req}
